\documentclass[10pt,a4paper]{amsart}
\usepackage[margin=19mm]{geometry}
\usepackage{amsmath,amssymb,mathtools}
\usepackage[scr=rsfs,cal=euler]{mathalfa}
\usepackage{xfrac}
\usepackage[unicode,hidelinks,bookmarks=false]{hyperref}
\newcommand{\ie}{i.e.\ }
\newcommand{\cf}{cf.\ }
\newcommand{\resp}{respectively\ }
\newcommand{\Subsection}{\subsection}
\providecommand{\nobreakdash}{\nobreak\hbox{-}}
\setlength{\emergencystretch}{2em}
\multlinegap0pt
\usepackage{amssymb}
\usepackage{bm}
\usepackage{algorithm}\usepackage{algpseudocode}
\newtheorem{theorem}{Theorem}
\title{A high-order, meshless, Lagrangian--Eulerian RBF-FD method for advection--diffusion--reaction on moving manifolds}
\author{Matthew Lowery, Grady B. Wright, Varun Shankar}
\date{Source: arXiv:2608.19384v1 (CC BY 4.0)}
\begin{document}
\maketitle

\noindent\emph{Source and licence.} A high-order, meshless, Lagrangian--Eulerian RBF-FD method for advection--diffusion--reaction on moving manifolds. Version arXiv:2608.19384v1 (CC BY 4.0), distributed by arXiv under the Creative Commons Attribution 4.0 International licence, http://creativecommons.org/licenses/by/4.0/, as recorded in the arXiv licence metadata for this exact version and shown on the abstract page https://arxiv.org/abs/2608.19384v1. Full licence notice: \texttt{licenses/arxiv-2608.19384-CC-BY-4.0.txt}.\par\smallskip

\noindent\textit{Editorial scope.} Counted unit: the article's theorem thm:approximate-normal-tangent-rbffd (source0.tex lines 1510--1544). The article's complete original proof of it is delivered with it (source0.tex lines 1546--1612, one \texttt{proof} environment of 67 lines). No mathematics is written, repaired or re-derived: the statement and the proof are the article's own lines, in the article's own notation, and no retained line is substituted or re-flowed.  Retained uncounted as the source-local context the proof needs: lines 1407--1436.  Nothing was omitted from any retained span, so the package carries no omission record.  No retained line contains a citation, so the wrapper carries no bibliography and no bibliography entry.  No retained line contains a figure environment, an image-inclusion command or drawing code, so no image is delivered and the package declares no asset dependency.  Editorial formatting only, in addition: an AMS \texttt{amsart} preamble that restates the article's own declarations and macros used by the retained lines, namely the environments \texttt{theorem} on separate counters, together with the standard packages those lines need.  The retained environments print in this document as Theorem 1 in order of appearance, whereas the article prints the same items with the numbering of its own full document.  The complete native source of the article as distributed in the arXiv e-print for this exact version is bundled unchanged as \texttt{sources/208085/source0.tex}.
\begin{theorem}[Convergence with exact normals]
\label{thm:exact-normal-tangent-rbffd}
Let \(\mathcal M\subset\mathbb R^3\) be a compact co-dimension one manifold that is closed and
\(C^{\xi+2}\).  Consider a sequence of point clouds \(X_h\subset\mathcal M\)
with \(h\to0\) (in the sense described above), and construct the tangent-plane RBF-FD operators using the
exact surface normals.  Suppose the formulas reproduce
\(\Pi_{\xi+1}(\mathbb R^2)\) and that there is a constant \(C_\Lambda\),
independent of \(h\), such that
\begin{equation}
  \Lambda_\nabla\leq C_\Lambda,
  \qquad
  \Lambda_\Delta\leq C_\Lambda
  \label{eq:theory-exact-uniform-stability}
\end{equation}
for every point cloud in the refinement sequence.  Then, for
\(u\in C^{\xi+2}(\mathcal M)\) and
\(U=(u(\bm x_1),\ldots,u(\bm x_N))^T\),
\begin{align}
  \max_{\alpha\in\{x,y,z\}}
  \|G_\alpha U-\partial_{\mathcal M,\alpha}u|_{X_h}\|_\infty
  &\leq C h^{\xi+1}\|u\|_{C^{\xi+2}(\mathcal M)},
  \label{eq:exact-normal-global-gradient}\\
  \|LU-\Delta_{\mathcal M}u|_{X_h}\|_\infty
  &\leq C h^\xi\|u\|_{C^{\xi+2}(\mathcal M)}.
  \label{eq:exact-normal-global-laplacian}
\end{align}
The constant \(C\) is independent of \(h\) and \(u\).  Hence the tangent-plane
gradient and Laplace--Beltrami approximations converge with orders
\(\xi+1\) and \(\xi\), respectively.
\end{theorem}

\begin{theorem}[Convergence with approximate normals]
\label{thm:approximate-normal-tangent-rbffd}
Let the assumptions of Theorem~\ref{thm:exact-normal-tangent-rbffd} hold, but
construct the tangent planes from approximate unit normals
\(\widetilde{\bm n}_i\).  Assume \(\varepsilon_h\leq\varepsilon_0<1\) for all
sufficiently small \(h\).  Also assume that there is a constant
\(C_\Lambda\), independent of \(h\), such that the approximate-plane weights
satisfy
\begin{equation}
  \max_{i,\alpha}
  \rho_i\|w_{i,\partial_{\mathcal M,\alpha}}\|_1
  \leq C_\Lambda,
  \qquad
  \max_i
  \rho_i^2\|w_{i,\Delta_{\mathcal M}}\|_1
  \leq C_\Lambda
  \label{eq:theory-approximate-uniform-stability}
\end{equation}
for every point cloud in the refinement sequence.  Then
\begin{align}
  \max_{\alpha\in\{x,y,z\}}
  \|G_\alpha U-\partial_{\mathcal M,\alpha}u|_{X_h}\|_\infty
  &\leq C\bigl(h^{\xi+1}+\varepsilon_h\bigr)
  \|u\|_{C^{\xi+2}(\mathcal M)},
  \label{eq:approximate-normal-global-gradient}\\
  \|LU-\Delta_{\mathcal M}u|_{X_h}\|_\infty
  &\leq C\bigl(h^\xi+\varepsilon_h\bigr)
  \|u\|_{C^{\xi+2}(\mathcal M)}.
  \label{eq:approximate-normal-global-laplacian}
\end{align}
In particular, if the normal approximation converges with order \(q_g\),
meaning $\varepsilon_h=O(h^{q_g})$,
then the surface gradient and Laplace--Beltrami errors are $O\!\left(h^{\min(\xi+1,q_g)}\right)$ and $O\!\left(h^{\min(\xi,q_g)}\right)$,
respectively.
\end{theorem}

\begin{proof}
Fix a stencil centered at \(\bm x_i\).  Since the approximate tangent plane
makes an angle \(O(\varepsilon_h)\) with the exact tangent plane, for
\(\varepsilon_h\) sufficiently small the surface can be written locally as a
graph over the approximate plane,
\[
  \widetilde\Psi_i(\bm y)
  =\bm x_i+\widetilde T_i\bm y
  +\widetilde{\bm n}_i\widetilde g_i(\bm y),
\]
where \(\widetilde T_i\) is an orthonormal basis of that plane.  At the stencil
center,
\[
  \|\nabla\widetilde g_i(0)\|_2\leq C\varepsilon_h,
\]
and the derivatives of \(\widetilde g_i\) needed below remain uniformly
bounded by the smoothness of \(\mathcal M\). Set \(\widetilde u_i=u\circ\widetilde\Psi_i\).  The induced metric is
\[
  \widetilde G_i
  =I+\nabla\widetilde g_i\nabla\widetilde g_i^T,
\]
so \(\widetilde G_i(0)^{-1}=I+O(\varepsilon_h^2)\).  The coordinate formula for
the surface gradient therefore gives
\[
  \left\|
    \nabla_{\mathcal M}u(\bm x_i)
    -\widetilde T_i\nabla\widetilde u_i(0)
  \right\|_2
  \leq C\varepsilon_h\|u\|_{C^1(\mathcal M)}.
\]
Similarly, the local-coordinate formula for \(\Delta_{\mathcal M}\) differs
from the Euclidean Laplacian by metric and first-derivative terms whose
coefficients are \(O(\varepsilon_h)\).  Hence
\[
  \left|
    \Delta_{\mathcal M}u(\bm x_i)
    -\Delta\widetilde u_i(0)
  \right|
  \leq C\varepsilon_h\|u\|_{C^2(\mathcal M)}.
\]

The Taylor argument from
Theorem~\ref{thm:exact-normal-tangent-rbffd}, applied in the approximate-plane
coordinates, gives for a derivative of order \(s\in\{1,2\}\),
\[
  \left|
    \widetilde{\mathcal D}\widetilde u_i(0)
    -\sum_{j=1}^{n_s}w_{ij}\widetilde u_i(\bm y_{i_j})
  \right|
  \leq
  C\Lambda_i\rho_i^{\xi+2-s}
  \|u\|_{C^{\xi+2}(\mathcal M)},
\]
where \(\Lambda_i=\rho_i^s\|w_i\|_1\).  By
\eqref{eq:theory-approximate-uniform-stability},
\(\Lambda_i\leq C_\Lambda\).  Combining this differentiation error with the
\(O(\varepsilon_h)\) geometric error yields
\[
  |E_i|
  \leq
  C\left(\rho_i^{\xi+2-s}+\varepsilon_h\right)
  \|u\|_{C^{\xi+2}(\mathcal M)}.
\]
Since \(\rho_i\leq h\), taking the maximum over all stencils proves
\eqref{eq:approximate-normal-global-gradient}--
\eqref{eq:approximate-normal-global-laplacian}.
\end{proof}

\end{document}
